\documentclass[letterpaper]{article}

\usepackage{graphicx}
\usepackage{amsfonts}
\usepackage{amsthm}
\usepackage{amsmath}
\usepackage{amssymb}
\usepackage{scribe}
\usepackage[margin=1.5in]{geometry}

\usepackage{braket}
\usepackage{xcolor}

\begin{document}

\makeheader{Kevin Sheng}{October 1, 2026}{3}{Using Fooling Sets}

\noindent
In the previous lecture, fooling sets were introduced.
Here we will prove some of their properties and see
that they can be a useful tool for proving lower bounds for some functions.

Recall the definition of a fooling set:
\begin{definition}
    A set $S \subseteq X \times Y$ is a \underline{\textbf{fooling set}}
    for $f: X \times Y \to \{0, 1\}$ if:
    \begin{enumerate}
        \item $f$ is constant on $S$.
        \item For all distinct pairs $(x, y), (x', y') \in S$, $f$ is \textbf{non-constant} on \\
              $\{(x, y), (x', y'), (x, y'), (x', y)\}$.
    \end{enumerate}
\end{definition}

\section{Challenge problem solution}

We first prove a challenge theorem from lecture 2:
\begin{problem}
Prove that for fooling set $\{(x_1, y_1), \cdots, (x_n, y_n)\}$,
$x_i \ne x_j$ and $y_i \ne y_j$ for any two distinct $1 \le i, j \le n$.
\end{problem}

\noindent
Since characteristic matrices in communication complexity are defined only up to
a permutation of the rows and columns, this formulation is equivalent to saying
that all fooling sets are subsets of generalized diagonals.

\begin{proof}
    We'll prove the contrapositive.
    Consider any $S=\{(x_1, y_1), \cdots, (x_n, y_n)\}$.
    Let there exist $1 \le i, j \le n$ s.t. $i \ne j$ and $x_i = x_j$.
    We can let $x$ take on the identical elements without loss of generality.

    By definition, $f$ must be constant on $\{(x_i, y_i), (x_j, y_j)\}$
    and at the same time be non-constant on $\{(x_i, y_i), (x_j, y_j), (x_i, y_j), (x_j, y_i)\}$.
    However, since $x_i=x_j$, this set is identical to $\{(x_i, y_i), (x_j, y_j)\}$
    after removing duplicate elements.
    Thus, $S$ can't be a fooling set.
\end{proof}

\section{Fooling sets as a lower bound}

A notable application of fooling sets is their use in proving
lower bounds on various functions' communication complexities.

To see this, we first need to establish the concept of a fooling set number.

\begin{definition}
    Given an $f: X \times Y \to \{0, 1\}$, its \underline{\textbf{zero fooling set number}}, $fs_0(f)$,
    is the largest fooling set that evaluates to all zeros.
    Its \underline{\textbf{one fooling set number}}, $fs_1(f)$, is defined analogously.
    The overall \underline{\textbf{fooling set number}} is $fs(f)=\max(fs_0(f), fs_1(f))$.
\end{definition}

With this definition in place, we can now prove the following theorem:

\begin{theorem}
    Given a function $f: X \times Y \to \{0, 1\}$, $\mathcal{D}(f) \ge \lceil\log fs(f)\rceil$.
\end{theorem}

\begin{proof}
    For convenience, let $c=\mathcal{D}(f)$.
    By Theorem 2.6, we can partition $M_f$ into $2^c$ disjoint $f$-monochromatic rectangles.

    Let $S$ be any fooling set of $f$.
    Notice that the definition of the fooling set combined with the locality property
    of rectangles makes it such that no two elements of $S$ can reside within the same rectangle.

    Thus, $2^c \ge |S|$ and by extension $c \ge \log |S|$.
    Since $c$'s a natural number, we can apply the ceiling operator to the RHS
    without any consequence, giving us $c \ge \lceil\log |S|\rceil$.
    Replacing $c$ with $\mathcal{D}(f)$ and using $f$'s largest fooling set
    for $S$ with size $fs(f)$ indeed gives the desired inequality.
\end{proof}

However, note that we made no distinction between the color of the fooling sets
and the color of the partitioned rectangles.
If we observe that a rectangle with only $0$s can contain elements from a zero
fooling set and similarly with a rectangle that has only $1$s,
we can actually prove a slightly stronger version:
\begin{theorem}
    Given a function $f: X \times Y \to \{0, 1\}$, $\mathcal{D}(f) \ge \lceil\log (fs_0(f)+fs_1(f))\rceil$.
\end{theorem}

\begin{proof}
    Let the variables be defined as they are in the previous proof.

    This time, let's have $S_0$ be a zero fooling set of $f$, and $S_1$ be a one fooling set of $f$.
    We need enough rectangles with zeros to contain every element of $S_0$,
    and we also need enough with ones for the elements of $S_1$.
    These two sets of rectangles are disjoint, so overall $2^c \ge |S_0|+|S_1|$.

    Once again, replacing $c$ with $\mathcal{D}(f)$ and using $f$'s largest
    one fooling set and zero fooling set yields the inequality we want.
\end{proof}

\begin{problem}
If $fs_0(f)=1$, can we bound or say anything about $fs_1(f)$?
\end{problem}

\section{Usage}

With this fooling set method, we can prove lower bounds on communication
complexity for a good number of problems for which we were previously unable to.

\begin{example}
    Consider the same $M_f$ near the end of Lecture 2:
    \[M_f=\begin{bmatrix}
            \color{red}{1} & 1 & 0              \\
            0              & 0 & \color{red}{1} \\
            1              & 1 & 0
        \end{bmatrix}\]
    It was established before that this has communication complexity of $2$ bits.

    We can construct a one fooling set of size $2$ if we take the highlighted red elements.
    It is possible to construct a fooling set of size $0$, but in reality
    all we need for the lower bound is the observation is that the matrix is
    non-constant, and therefore $fs_0(f) \ge 1$.

    We can then lower bound $\mathcal{D}(f)$ like so:
    \begin{align*}
        \mathcal{D}(f)
         & \ge \lceil\log (fs_0(f)+fs_1(f))\rceil \\
         & \ge \lceil\log (1 + 2)\rceil           \\
         & = 2
    \end{align*}
    which shows that our $2$-bit protocol was the best we could've done.
\end{example}

\begin{example}
    In the first lecture, $EQ_n: \{0, 1\}^n \times \{0, 1\}^n \to \{0, 1\}$ is defined as
    \[EQ_n(x, y)=\begin{cases}
            1 & x = y            \\
            0 & \text{otherwise}
        \end{cases}\]
    We initially came up with a protocol that needed $n+1$ bits, which seemed inefficient.
    As we will see, this protocol is actually the best we can do deterministically.

    As long as we give the same order to the rows and columns in the characteristic matrix,
    $M_{EQ_n}$ looks like this $2^n \times 2^n$ matrix:
    \[
        \left[
            \begin{array}{ccccc}
                1                                             \\
                 & 1             &        & \text{\huge0}     \\
                 &               & \ddots                     \\
                 & \text{\huge0} &        & 1                 \\
                 &               &        &               & 1
            \end{array}
            \right]
    \]
    Here we can actually construct a one fooling set of maximum size
    by taking the diagonal. Any crosspair gives us $0$s on both ends
    of the formed rectangle, so we can include all the terms.
    Once again noticing that $fs_0(EQ_n) \ge 1$, we have
    \[\mathcal{D}(EQ_n) \ge \lceil\log\left(2^n+1\right)\rceil = n+1\]
    which indeed proves that the naive protocol was as good as we could do.
\end{example}

\begin{problem}
What's the exact value of $fs_0(EQ_n)$?
\end{problem}

\begin{example}
    If we map $\{0, 1\}^n$ to the natural numbers using some $g: \{0, 1\}^n \to \mathbb{N}$,
    we can define $GT_n: \{0, 1\}^n \times \{0, 1\}^n \to \{0, 1\}$ like so:
    \[GT_n(x, y)=\begin{cases}
            1 & g(x) \ge g(y)    \\
            0 & \text{otherwise}
        \end{cases}\]
    Despite the name being the abbreviation of "greater than", this function
    actually just checks for whether one number is greater than \textit{or equal to} the other.

    Here, the characteristic matrix takes the following form:
    \[
        \left[
            \begin{array}{ccccc}
                1                                             \\
                 & 1             &        & \text{\huge0}     \\
                 &               & \ddots                     \\
                 & \text{\huge1} &        & 1                 \\
                 &               &        &               & 1
            \end{array}
            \right]
    \]
    (Inputs to $x$ are down along the rows, while inputs to $y$ are across the columns.)

    Like with $EQ_n$, we can create a one fooling set of maximum size by taking the diagonal.
    While the lower triangle is now all $1$s, the upper triangle provides all the $0$s
    we need to make any crosspair non-constant.
    With that, we have our lower bound on $GT_n$:
    \[\mathcal{D}(GT_n) \ge \lceil\log\left(2^n+1\right)\rceil = n+1\]
\end{example}

\begin{example}
    Another canonical problem in communication complexity is the disjointness problem
    $DISJ_n: \mathcal{P}(\{1, 2, \cdots, n\}) \times \mathcal{P}(\{1, 2, \cdots, n\}) \to \{0, 1\}$:
    \[DISJ_n(x, y)=\begin{cases}
            1 & x \cap y = \varnothing \\
            0 & \text{otherwise}
        \end{cases}\]
    A real-life example of this could come up in meeting scheduling, where
    Alice and Bob each have their own sets of days on which they're free
    and the two need to find out whether a meeting is even possible in the first place.

    Simply creating the outline of the characteristic matrix and pattern matching
    won't get us very far now.
    Instead, let's try and think about how we might analytically construct
    a maximal one fooling set.

    Any pair in this set should have two general properties:
    \begin{enumerate}
        \item Consist of two disjoint sets (obviously).
        \item Cover as much of the universe as possible, so that
        when we swap out one of the pairs for some other arbitrary set
        the likelihood of a common element with the preserved set is maximized.
    \end{enumerate}

    This leads to the following candidate for a fooling set:
    \[S=\left\{\left(A, \overline{A}\right) \mid A \subseteq \{1, 2, \cdots, n\}\right\}\]
    Notice that order matters here.
    The set $\{1\}$ occurs as both the first element and second element in a pair.
    \begin{claim}
        $S$ is a fooling set.
    \end{claim}
    \begin{proof}
        Since a set and its complement cannot share any elements, every pair in $S$ is disjoint.

        To show that crosspairs aren't monochromatic, take any two distinct sets $A \ne B$.
        It STP either $DISJ_n(A, \overline{B})=0$ or $DISJ_n(B, \overline{A})=0$.

        WLOG assume $\exists x \in A: x \notin B$.
        This implies $x \in \overline{B} \implies x \in A \cap \overline{B}$,
        which means that $DISJ_n(A, \overline{B})=0$.
    \end{proof}
\end{example}

\begin{problem}
Show that one can rearrange the rows and columns in $M_{DISJ_n}$ to
make it take the below form:
\[
    \left[
        \begin{array}{ccccc}
            1                                             \\
             & 1             &        & \text{\huge*}     \\
             &               & \ddots                     \\
             & \text{\huge0} &        & 1                 \\
             &               &        &               & 1
        \end{array}
        \right]
\]
where the $*$ on the top right means we don't care what happens in the top right triangle.
This leads to another proof of the communication complexity lower bound,
as one can take the diagonal to get a maximal one fooling set.
\end{problem}

Many problems are like the ones just discussed in that they have
very high lower bounds on the deterministic complexity.
However, they can still fall with the application of more advanced techniques,
such as probabilistic or multi-party communication.
This is with the sole exception of one $IP_n$, of course, which we discuss below:

\begin{example}
    $IP_n: \{0, 1\}^n \times \{0, 1\}^n \to \{0, 1\}$ is defined as follows:
    \[IP_n(x, y)=\bigoplus_{i=1}^n x_i \land y_i\]
    where $\land$, despite being the XOR operator in Python and many other
    languages, indicates the bitwise AND operator here.

    This function can be thought of as taking the inner product of two vectors
    that are in $\mathbb{F}_2$, the simplest field with only two elements ($0$ and $1$).
    Expanding more on this line of thought, let's try and use it to create
    a fooling set.
    
    The idea is to construct vectors have a dot product of $1$ with themselves
    but a $0$ dot product with others.
    In other words, we want $v_1, \cdots, v_k \in \mathbb{F}_2^n$ s.t.:
    \begin{enumerate}
        \item $\braket{v_i, v_i}=1$
        \item $\braket{v_i, v_j}=0\ \forall i \ne j$
    \end{enumerate}
    These conditions allow us to declare
    $\{(v_1, v_1), (v_2, v_2), \cdots, (v_k, v_k)\}$ a fooling set.

    However, these constraints, which seem eerily similar to the conditions
    for an orthonormal basis, severely limit how big $k$ can get:
    \begin{claim}
        Vectors $v_1$ through $v_k$ are linearly independent.
    \end{claim}
    \begin{proof}
        Suppose there existed coefficients $\lambda_1, \cdots, \lambda_k$
        s.t. $\sum_{i=1}^{k} \lambda_i v_i = \mathbf{0}$.
        Then for any $1 \le j \le k$ that we choose, we have
        \[\Braket{\mathbf{0}, v_j}=0 \implies \Braket{\sum_{i=1}^{k} \lambda_i v_i, v_j}=0\]
        By the properties of the inner product, we can take the summation
        outside to obtain $\sum_{i=1}^{k} \lambda_i\Braket{v_i, v_j}=0$.
        For all terms in the summation where $i \ne j$, $\Braket{v_i, v_j}=0$.
        This leaves $\Braket{v_j, v_j}$ as the only nonzero dot product,
        which forces $\lambda_j=0$.
        However, the choice of $j$ at the beginning was arbitrary,
        which means that $\lambda_i=0$ for all $i$ from $1$ through $k$.
    \end{proof}
    This linear independence constraint bounds our fooling set construction
    by $n$, which gives us a pitiful lower bound of
    \[\mathcal{D}(f) \ge \lceil\log(n+1)\rceil\]
\end{example}

This could have been because we forced the two elements of a fooling pair to be the same.
Could we have done better by removing this constraint?

\begin{problem}
Show that the fooling set number of $IP_n$ is bounded by
some expression that is polynomial in $n$.
That is,
\[fs(IP_n) \le \text{poly}(n)\]
\end{problem}

\section{A typical fooling set}

Now we shift our attention to what a "typical" maximal fooling set would look like.

Unfortunately, we have the following theorem:
\begin{theorem}
    For a random $f: \{0, 1\}^n \times \{0, 1\}^n \to \{0, 1\}$,
    \[\mathbb{P}[fs(f) \ge 10n] \le \left(\frac{2}{3}\right)^{n^2}\]
\end{theorem}
In other words, the chance we can obtain a fooling set that
gives us a logarithmic lower bound plus some constant factor
vanishes at a rate faster than standard exponential decay.

\begin{proof}
    Instead of thinking about all these functions as random maps,
    let's instead think about them as random bit matrices.
    The idea is to consider all fooling set candidates of size $10n$,
    estimate their probability of being a valid fooling set,
    and then loosely bound $\mathbb{P}[fs(f) \ge 10n]$ from there.

    First fix some arbitrary fooling set candidate of size $10n$.
    Two things have to be true:
    \begin{enumerate}
        \item $f$ evaluates to the same element on every argument pair in this set.
        \item Every crosspair gives at least one element that's different from the diagonal.
    \end{enumerate}
    The chance of the first occurring is $\frac{1}{2^{10n-1}}$,
    since there's $2$ possible values and the chance that all elements
    along the diagonal evaluate to that one value is $\frac{1}{2^{10n}}$.

    Evaluating the chance of the second is a little trickier,
    but a key observation is that crosspairs are both independent of the main
    diagonal and each other.
    If we need $f$ to be nonconstant across the crosspair of $(a, b)$ and $(c, d)$,
    we need at least one different element in either $(a, d)$ or $(c, b)$.
    Neither of these pairs are in the actual fooling set candidate,
    and they also aren't looked at by any other crosspair in the fooling set candidate.
    Thus, the total chance of the second condition being met is
    $\left(\frac{3}{4}\right)^{\tbinom{10n}{2}}$.

    Combining these probabilities, we see that the chance of any random
    fooling set candidate actually being a fooling set is
    $\frac{1}{2^{10n-1}} \cdot \left(\frac{3}{4}\right)^{\tbinom{10n}{2}}$.

    In total, there's $\binom{2^n}{10}$ ways to select the first elements
    of the fooling pairs and the same number of ways to select the second elements.
    Thus, we can use the union bound to see that as long as $n$ is large enough,
    \begin{align*}
        \mathbb{P}[fs(f) \ge 10n]
        &\le \binom{2^n}{10}^2\left(\frac{1}{2^{10n-1}} \cdot \left(\frac{3}{4}\right)^{\tbinom{10n}{2}}\right) \\
        &\le \left(\left(2^n\right)^{10}\right)^2\left(\frac{1}{2^{10n-1}} \cdot \left(\frac{3}{4}\right)^{\tbinom{10n}{2}}\right) \\
        &\le 2^{20n} \cdot \frac{1}{2^{10n-1}} \cdot \left(\frac{3}{4}\right)^{50n^2-5n} \\
        &\le 2^{10n+1} \cdot \left(\frac{3}{4}\right)^{50n^2} \\
        &\le 2^{10n+1} \cdot \left(\frac{2}{3}\right)^{25n^2} \\
        &\le \left(\frac{2}{3}\right)^{n^2}
    \end{align*}
\end{proof}

\end{document}
